% Copyright (c) 2026 Eugene Joseph M. Ragasa.
% SPDX-License-Identifier: Apache-2.0
% Author: Eugene Joseph M. Ragasa
% Formal proof completed: 2026-10-20
\documentclass[11pt]{article}
\usepackage[T1]{fontenc}\usepackage{amsmath,amssymb}\usepackage[margin=1in]{geometry}\usepackage[hidelinks]{hyperref}
\title{Proof Dossier: Joint-Enclosure Check Implies Validity}\author{Eugene Joseph M. Ragasa}\date{Formal proof completed: 2026-10-20}
\begin{document}\maketitle
\section{Identity and status}
ID \texttt{BG-CERT-THM-013}; status \texttt{proved}. Lean declaration:
\par\noindent{\tiny\path{ProjectKoios.BandGauge.Certificates.}}\\
{\tiny\path{valid_of_checkRankTwoSeparationBoxCertificate_eq_true}}.
\section{Exact statement}
Boolean acceptance of a rational joint $(\kappa,\mu)$ enclosure certificate implies every
encoded validity clause for positivity, endpoint order, witness, radii, containment, and
strict separation.
\section{Proof and boundary}
The checker is the decision procedure for the explicit rational conjunction. The proof
uses \path{of_decide_eq_true} and repackages the clauses in the named validity
structure. It provides no parser, source identity, runtime artifact, or scientific
validation.

Kernel axioms are \texttt{propext}, \texttt{Classical.choice}, and
\texttt{Quot.sound}; no admissions. Checked with Lean \texttt{4.35.0} and Mathlib
revision \texttt{56a005855ea17e0b0958f73ec801ae140b0a1679}. AI-assisted and
kernel-checked; no independent expert review recorded.
\end{document}
